\section{Rational functions on odd grids}\label{sec:construction}

\subsection{Zeros, poles and differentiation}

Fix a positive odd integer $m$, a count $B\ge1$ and an integer
$r\ge1$. These parameters remain fixed as the positive odd
integer $n$ tends to infinity. Put
\[
 k=m-1,\qquad M=2B-1,\qquad L=2r+1.
\]
The $b$th form will be an alternating sum over the grid $j/b$
with $j$ odd, for $b=1,3,\ldots,M$. We need zeros on these
grids both to remove initial summands and to control the
denominators of the rational constants.
We place the poles at $0,-2,\ldots,-2n$ and the grid zeros in
$(-(2r+2)n,2rn)$. This interval is centered at $-n$ and has
length $2Ln$; $r$ controls how far it extends beyond the poles.

The reduced fraction classes modulo $2$ are
\[
 \Theta=\left\{\frac vd:
 1\le d\le M\text{ odd},\ 1\le v<2d\text{ odd},\ (v,d)=1\right\}.
\]
There are $\varphi(d)$ classes for each $d$, including the single
class $1$ for $d=1$. Thus $|\Theta|=D_B$; write $D=D_B$.
The integer $D$ is odd, since $\varphi(d)$ is even for odd $d>1$.
Each class contributes $Ln$ points to the zero interval, spaced
by $2$. The polynomial $P_n$ below therefore has $DLn$ simple zeros.
Define
\begin{equation}\label{eq:K}
 A_d=(2d)^L\prod_{p\mid d}p^{\lceil L/(p-1)\rceil},
 \qquad K=\prod_{\theta\in\Theta}A_{\operatorname{den}(\theta)}.
\end{equation}
The product over primes supplies the factorial divisibility
proved in Section~\ref{sec:arithmetic}.

Choose an even pole order $a>mDL$. Set
\begin{align}
 P_n(t)&=\prod_{\theta\in\Theta}\prod_{j=0}^{Ln-1}
                      (t-2rn+2j+\theta),\label{eq:P}\\
 \kappa_n&=2^{an}K^{mn}(n!)^{a-mDL},
 &R_n(t)&=\frac{\kappa_nP_n(t)^m}{\prod_{h=0}^n(t+2h)^a}.
 \label{eq:R}
\end{align}
The factorial in $\kappa_n$ balances the growth of the products;
the powers of $2$ and $K$ give
the coefficient divisibility proved in Section~\ref{sec:arithmetic}.
None of these grid zeros is a pole.
The map $\theta\mapsto2-\theta$ preserves $\Theta$, so the root
set is invariant under $t\mapsto-2n-t$. Since $DLn$ and $m$
are odd and $a$ is even,
\[
 P_n(-2n-t)=-P_n(t),\qquad R_n(-2n-t)=-R_n(t).
\]

Multiplication by an integer polynomial will let us adjust the
rational constants of the forms, particularly at the seed index
$n=1$, while retaining all grid zeros. We choose the polynomial
even about $-n$ to preserve reflection.
Section~\ref{sec:gram} constructs such a polynomial $W_n(t)$,
homogeneous in $n,t$ of total degree
\begin{equation}\label{eq:E}
 E=m(D-1).
\end{equation}
Homogeneity will control its scaling at prime indices.
Define
\begin{equation}\label{eq:T-V}
 T_n(t)=W_n(t)R_n(t),\qquad V_n(t)=\frac1{k!}T_n^{(k)}(t).
\end{equation}
Both functions are antisymmetric about $-n$, because $k$ is
even. Every grid zero has multiplicity at least $m=k+1$ in
$T_n$, so $V_n$ still vanishes there. The degree of $T_n$ at
infinity is at most
\[
 -(a-mDL)n-a+E.
\]
Since $a>mDL$ and $L\ge3$, we have $a>E$; the functions are
proper and the sums below converge absolutely.

\subsection{The linear forms}

Write the partial fractions before and after differentiation as
\begin{align}
 T_n(t)&=\sum_{h=0}^n\sum_{i=1}^a
                     \frac{c_{i,h}(n)}{(t+2h)^i},\notag\\
 V_n(t)&=\sum_{h=0}^n\sum_{\nu=m}^{a+k}
                     \frac{A_{\nu,h}(n)}{(t+2h)^\nu}.
 \label{eq:partials}
\end{align}
Differentiation relates their coefficients by
\[
 A_{i+k,h}(n)=\binom{i+k-1}{k}c_{i,h}(n),
 \qquad 1\le i\le a.
\]
The coefficients $A_{\nu,h}$ are zero outside $m\le\nu\le a+k$.
Put
\[
 \rho_\nu(n)=\sum_{h=0}^n(-1)^h A_{\nu,h}(n).
\]
Antisymmetry gives
$A_{\nu,n-h}=(-1)^{\nu+1}A_{\nu,h}$. Since $n$ is odd,
$\rho_\nu=(-1)^\nu\rho_\nu$, and every odd coefficient vanishes.

For odd integers $j$ write $\ch(j)=(-1)^{(j-1)/2}$.
For each odd $b\le M$, put
\[
 F_{n,b}=\sum_{\substack{j\ge1\\j\text{ odd}}}\ch(j)V_n(j/b).
\]
In the term from the pole $-2h$, shift the character index by
$2bh$. Since $b$ is odd, the sign changes by $(-1)^h$.
The omitted initial terms give
\begin{align}
 F_{n,b}&=q_b(n)+
   \sum_{\substack{m+1\le\nu\le a+m-1\\\nu\text{ even}}}
             b^\nu\rho_\nu(n)\beta(\nu),\label{eq:forms}\\
 q_b(n)&=-\sum_{h=0}^n(-1)^h\sum_{\nu=m}^{a+k}A_{\nu,h}(n)b^\nu
                  \sum_{\substack{1\le j<2bh\\j\text{ odd}}}
                             \frac{\ch(j)}{j^\nu}.
 \label{eq:q}
\end{align}
At $m=1$, the exponent-one partial-fraction series converge
conditionally. Only finitely many series are combined, and
their total $\beta(1)$ coefficient is zero; no infinite
rearrangement is involved.

\subsection{A rational criterion for the endpoints}

The positive rational $u$ below bounds the location of the largest
summands, measured beyond $2rn$ in units of $n$. The quantity $U$
bounds the square of the exponential rate after clearing
denominators.

\begin{proposition}\label{prop:criterion}
With the parameters above, let $e_*>\ep$ be rational. Suppose
a positive rational $u$ satisfies
\begin{align}
 f(u)&=\left(\frac{u+2L}{u}\right)^{mD}
       \left(\frac{u+2r}{u+2r+2}\right)^a<1,\label{eq:f}\\
 U&=e_*^{2(a+k)}K^{2m}(u+2L)^{2mDL}2^{2a}
       \left(\frac{(2r)^{2r}}{(2r+2)^{2r+2}}\right)^a<1.
 \label{eq:U}
\end{align}
Then at least $B$ of
$\beta(m+1),\beta(m+3),\ldots,\beta(a+m-1)$ are irrational.
\end{proposition}

The exponent $a+k$ pays for denominators after differentiation.
To finish the criterion, we also need a seed constant that survives
the elimination of the possible irrational values. We will choose
one multiplier $W_n$ shared by all the forms so that
\begin{equation}\label{eq:seed-target}
 (q_1(1),q_3(1),\ldots,q_M(1))=\eta e_B,
 \qquad \eta\in\Q\setminus\{0\}.
\end{equation}
Here $e_B=(0,\ldots,0,1)$ has $B$ coordinates.
The elimination can be performed with nonzero weight on the last
form, so this seed vector retains a nonzero constant.
